\section{Related Work}\label{sec:related}

\subsection{Video Diffusion Models}

Diffusion models~\cite{first_diffusion, ddpm, ddim} have demonstrated unprecedented power in video generation. Video Diffusion Models (VDMs)~\cite{vdm} are broadly categorized into Text-to-Video (T2V) and Image-to-Video (I2V) frameworks, aiming to generate video samples from text prompts or image prompts. T2V generation~\cite{align_your_latents, videocrafter_1, dreamix, make_a_video, SVD, phenaki, tune_a_video, cogvideo, cogvideox, khachatryan2023, ge2023, animatediff} has been extensively studied in recent years, introducing text descriptions to control the content of video generation semantically. Previous works~\cite{align_your_latents, lvdm, modelscope, lavie, magicvideo, animatediff} incorporate temporal layers into large pretrained text-to-image (T2I) diffusion models~\cite{rombach2022}. Subsequent studies~\cite{align_your_latents, SVD, cogvideox, opensora, opensora_plan} have expanded T2V capabilities by utilizing large text-video pairs, achieving improved results. Building upon T2V, I2V synthesis~\cite{i2vgen_xl, SVD, cogvideox, opensora, videocrafter2, dynamicrafter, liu2024dynamic} has also been widely explored. Given a still image, I2V aims to animate it into a video clip that retains all visual content from the image and exhibits naturally suggested dynamics. Many recent works, such as SVD~\cite{SVD}, VideoCrafter2~\cite{videocrafter2} and CogVideoX~\cite{cogvideox} support both T2V and I2V simultaneously.


Despite producing high-quality videos, these models rely on text or image prompts, limiting fine-grained control and potentially leading to actions misaligned with user intentions. For precise control, some works~\cite{hu2024animate,huang2024lvcd,wang2024framer,zhu2024zero,xu2024magicanimate,ma2024follow,qin2023dancing,zhang2024mimicmotion,xing2023make,zhu2024champ,jeong2023tpos,zhang2025motion} employ multimodal video sequences as conditions, such as pose~\cite{zhu2024zero,xu2024magicanimate,hu2024animate,ma2024follow,qin2023dancing,zhang2024mimicmotion}, depth~\cite{xing2023make,he2023animate,zhu2024champ}, or sound~\cite{jeong2023tpos,lee2023aadiff,liu2023generative,tian2024emo}, treating video generation as a video translation task. Although these models achieve precise control, they require per-frame dense control signals, which makes them cumbersome and not user-friendly in real-world applications. Therefore, simpler yet precise control mechanisms are needed. Trajectory-based control offers an effective method for manipulating video generation, combining simplicity with precision.


\subsection{Trajectory Control in Video Generation}

Controllable editing has gained advancements in the field of image editing due to its precise control information~\cite{EpsteinJPEH23,MuGZSVWP24,YenphraphaiPLPX24,BhatMW24}. For video synthesis, trajectory-controlled generation has recently gained popularity due to its ability to achieve precise motion control. Early works~\cite{hao2018controllable, ardino2021click, blattmann2021understanding, ipoke} employed recurrent neural networks or optical flow to guide motion. Methods like TrailBlazer~\cite{trailblazer} utilize bounding boxes to direct subject motion in video generation. MotionCtrl~\cite{Motionctrl} encodes trajectory coordinates into dense vector maps, and DragNUWA~\cite{DragNUWA} transforms sparse strokes into dense flow spaces; both use these representations as guidance signals. Tora~\cite{Tora} employs a motion variational autoencoder~\cite{vae} to embed trajectory vectors into the latent space, preserving motion information across frames.

Although these methods facilitate trajectory control, they often lack semantic understanding of entities, making control over video generation less refined. To address this issue, DragAnything~\cite{DragAnything} combines entity representation extraction with a 2D Gaussian representation to achieve entity-level controllable video generation. TrackGo~\cite{TrackGo} uses user-provided free-form masks and arrows to define target regions and movement trajectories, serving as precise blueprints for video generation. However, all these methods consider 2D trajectories in image space, leading to ambiguities in real 3D environments. The recent 3D-TrajMaster~\cite{3dtrajmaster} manipulates multi-entity 3D motions with user-desired 6DoF pose sequences of entities for video generation. In this paper, we introduce an innovative control signal representation that combines depth information with K-means clustered points from object masks in video, achieving accurate entity-level and 3D trajectory control.